\section{Extended Classification of Unliftable Types} \label{sec:classification}

\subsection{Types with No Binary Operation At All}

\ps{See my point above about inducing it from the transition function.}

Consider a type $T$ that carries no natural merge structure whatsoever. The canonical example is:
\[
T = \{0, 1\}^* / \sim
\]
the set of bitstrings modulo some equivalence that is not a congruence with respect to any commutative operation. More concretely:

\begin{corollary}[Last-Write-Wins Register]
Let $T$ be any set, with the intended sequential semantics being that the most recent write wins. The ``operation'' is:
\[
a \oplus b = b \quad \text{(take the rightmost value)}
\]
This is not commutative since $a \oplus b = b \neq a = b \oplus a$ in general. Therefore no $\R$-lifting exists.
\end{corollary}

However, this is not because the type is unliftable in principle. It is because the sequential specification itself is inherently order-dependent, and no CRDT can implement it without additional metadata. The standard resolution is to enrich $T$ with timestamps, yielding a different sequential type $(T \times \text{Timestamp}, \oplus_{\text{LWW}}, \bot)$ that is liftable. The original $T$ is unliftable but only because it is under-specified. \qed


\subsection{Irreparably Unliftable Types}

The deepest obstruction arises from types whose semantics are inherently non-monotone. A function $f : T \to T$ is monotone with respect to a partial order if:
\[
a \leq b \implies f(a) \leq f(b)
\]
CRDTs are monotone by construction --- state only moves upward in the semilattice order. Therefore:

\begin{theorem}[Monotonicity Obstruction]
If the sequential specification of $(T, \oplus, e)$ requires any operation that is non-monotone with respect to every partial order on $T$, then $T$ admits no $\R$-lifting.
\end{theorem}

The canonical example is:
\[
T = (\mathbb{Z}, \min, +\infty)
\]
considered not as an output type but as a type supporting decrement without bound. A counter that can decrease arbitrarily has no monotone CRDT representation, because any decrease would require moving downward in the lattice order, contradicting the SEC requirement that state transitions are monotone.

More precisely:

\begin{theorem}[Non-liftability of unbounded decrement]
Let $T = (\mathbb{Z}, +, 0)$ with the operation set including both increment and decrement as primitive sequential operations rather than as derived from a PN-Counter lifting. Then $T$ admits no $\R$-lifting that is monotone and satisfies SEC.
\end{theorem}

\begin{proof}
Suppose $T\R$ exists. Since SEC requires that the state lattice is a join-semilattice, the state space must have a partial order $\leq$ such that every update is monotone increasing. But the decrement operation $d : T\R \to T\R$ must satisfy:
\[
\get(d(a)) = \get(a) - 1 < \get(a)
\]

So $d$ strictly decreases the observable value. For $d$ to be monotone, we need $d(a) \geq a$ in the lattice order, yet $\get$ must be order-preserving, so $\get(d(a)) \geq \get(a)$, contradicting $\get(d(a)) = \get(a) - 1$. Therefore no such lifting exists. \qed
\end{proof}

\medskip
\noindent\textbf{Remark.} The PN-Counter resolves this not by lifting $(\mathbb{Z}, +, 0)$ directly, but by changing the type to a pair of monotone counters $(\mathbb{N} \times \mathbb{N}, \oplus_{\text{PN}}, (0,0))$ whose $\get$ computes the difference. The integer type itself remains unliftable.

\begin{table*}[h]
    \centering
         \begin{tabular}{|l|c|l|}
\hline
\textbf{Type class} & \textbf{Liftable?} & \textbf{Reason} \\
\hline
Commutative monoids & \checkmark Yes & Constructive lifting exists \\
$(\mathbb{N}, +, 0)$, $(\mathbb{B}, \lor, \bot)$ & & (Theorem~1) \\
\hline
Non-commutative monoids & $\times$ No & Commutativity of $\sqcup$ forces \\
$(\Sigma^*, \cdot, \varepsilon)$ & & commutativity of $\oplus$ \\
\hline
Non-associative types & $\times$ No & Associativity of $\sqcup$ forces \\
$(\mathbb{Z}, -, 0)$ & & associativity of $\oplus$ \\
\hline
Non-monotone types & $\times$ No & SEC monotonicity is violated \\
$(\mathbb{Z}, +, 0)$ with decrement & & \\
\hline
Under-specified types & $\times$ No directly & Liftable only after adding \\
(LWW register) & \checkmark with enrichment & metadata to the type \\
\hline
    \end{tabular}
    \caption{classification of types with respect to liftability}
    \label{tab:liftabletypeclassification}
\end{table*}
